\section{Roadmap de Implementação}

\subsection{Visão Geral}

A implementação da solução será realizada de forma incremental, priorizando a entrega rápida de valor e a redução de riscos no processo de migração das integrações.

O roadmap foi estruturado em fases, permitindo validar a abordagem proposta em um escopo controlado antes de expandir sua aplicação para outras locadoras e fluxos.

\subsection{Estratégia de Execução}

Para garantir foco e eficiência na execução, propõe-se a criação de uma squad temporária dedicada, composta pelos desenvolvedores atualmente envolvidos no contexto de integrações.

Essa squad terá como objetivo atuar exclusivamente na evolução do processo de migração, incluindo:

\begin{itemize}
    \item Desenvolvimento do módulo de validação (Shadow Mode + Golden Master);
    \item Suporte à migração de integrações existentes;
    \item Evolução incremental da solução conforme aprendizado.
    \item Manutenção da esteira presente, culminando na continuação das entregas de tarefas extras sem necessidade de interromper quaisquer atividades.
\end{itemize}

A proposta não implica aumento de capacidade, mas sim uma reorganização temporária do time, visando reduzir o contexto paralelo de trabalho e acelerar a entrega.

Após a estabilização do sistema e conclusão das principais migrações, a squad poderá ser dissolvida, retornando ao modelo organizacional atual.

\subsection{Fases de Implementação}

\vspace{0.5em}
Fase 1 - Estruturação inicial
\begin{itemize}
    \item Definição da arquitetura do módulo de validação;
    \item Implementação básica do Shadow Mode;
    \item Captura das respostas do Rest e do Integrador;
    \item Implementação inicial do comparador (Golden Master);
    \item Persistência básica dos resultados.
\end{itemize}

Objetivo: validar a viabilidade técnica da solução.

\vspace{0.5em}
Fase 2 — Validação controlada
\begin{itemize}
    \item Aplicação do MVP em uma locadora/driver específico;
    \item Execução com volume controlado de requisições (amostragem);
    \item Identificação e análise de divergências;
    \item Ajuste das regras de comparação.
\end{itemize}

Objetivo: validar a abordagem em cenário real.

\vspace{0.5em}
Fase 3 — Expansão gradual
\begin{itemize}
    \item Ampliação para novos fluxos (ex: busca, reservar, cancelar, visualizar);
    \item Inclusão de novas locadoras no processo de validação;
    \item Evolução da classificação de divergências;
    \item Melhorias na persistência e estrutura dos dados.
\end{itemize}

Objetivo: escalar a solução de forma controlada.

\vspace{0.5em}
Fase 4 — Suporte à migração
\begin{itemize}
    \item Uso do sistema como base para decisões de migração;
    \item Execução de validações contínuas até estabilização;
    \item Apoio à conclusão da migração das integrações restantes.
\end{itemize}

Objetivo: garantir migração segura e confiável.

\vspace{0.5em}
Fase 5 — Evolução para observabilidade contínua
\begin{itemize}
    \item Estruturação do dashboard com métricas consolidadas;
    \item Evolução da solução para uso contínuo após migração;
    \item Integração com outras iniciativas de observabilidade.
\end{itemize}

Objetivo: transformar o sistema em um ativo permanente.